\section{SAT and exact counting below four gates per input}
\label{sec:algorithm-overview}

The algorithmic entry point is a restricted-instance improvement over
exhaustive search. Let $C$ be a single-output Boolean circuit with $u$
input bits and $s$ internal gates. Work over the full binary basis $B_2$:
every Boolean function of at most two inputs is a unit-cost gate,
constants and wiring are free, and fan-out is unrestricted.

The construction exploits two complementary measures: the number of
used inputs and the size of the boundary needed to maintain consistency.
The worked example establishes the local mechanism for summarizing a
region by conditional counts. The general argument separates exactness
of these summaries from the existence of an efficiently constructed
narrow processing order. The counting invariant establishes correctness;
the graph theorem supplies the resource bound.

\paragraph{Headline guarantee.}
For every fixed $\varepsilon>0$, the construction counts the satisfying
assignments of $C$, and therefore decides its satisfiability, in time
\begin{equation}\label{eq:headline-sat}
 \boxed{\quad
 \poly_\varepsilon(s+u+1)\,
 2^{(1/4+\varepsilon)(s+1)}.
 \quad}
\end{equation}
The space bound allows the same exponential factor. Consequently, for
each fixed $0<\gamma<4$, circuits with $s\le(4-\gamma)u$ admit an
algorithm running in $\poly_\gamma(u)2^{(1-\eta)u}$ for some
$\eta>0$ depending only on $\gamma$. This beats the $2^u$ enumeration
bound by an exponential factor. The restriction concerns the gate
budget in $B_2$, not the number of clauses in a CNF representation.

For example, take $\varepsilon=\gamma/[8(4-\gamma)]$ in
\eqref{eq:headline-sat}. Then
\[
 (1/4+\varepsilon)(s+1)
 \le (1-\gamma/8)u+(1/4+\varepsilon),
\]
so $\eta=\gamma/8$ is an explicit choice. The additive constant is
absorbed into the prefactor. The positive slack is fixed before the
input size grows; no exact endpoint exponent with a fixed polynomial
prefactor is asserted.

\paragraph{Why the running time improves.}
Normalize the output cone without changing its function, leaving $b$
used inputs. Each removed input contributes a factor of two to the
final count. Encode the circuit's gate equations, consistent uses of
every input and gate value, and the requirement that its output is one
as a network of small tables. Every satisfying input assignment has
exactly one extension to the internal values. The network's cycle rank
$\kappa$ satisfies
\[
 \kappa\le s+1-b.
\]
Split high-degree equality tables into trees, then contract low-degree
parts, loops, and parallel edges with exact table identities. A
nonempty residual simple cubic graph has at most $2(\kappa-1)$
vertices. The constructive Fomin--H\o ie theorem
\citep{FominHoie2006}, together with the width conversion proved below,
gives a vertex order with at most $(1/3+\delta)\kappa+O_\delta(1)$
edges crossing any frontier. Summing over the frontier assignments
evaluates the network using integer addition and multiplication.
Intermediate counts have polynomial bit length.

Choosing between this contraction and enumeration of the $b$ used
inputs gives the sharper bound
\begin{equation}\label{eq:headline-sensitive}
 \poly_\delta(s+u+1)\,
 2^{\min\{b,(1/3+\delta)(s+1-b)\}},
 \qquad \delta>0\text{ fixed}.
\end{equation}
Writing $\alpha=1/3+\delta$, the two exponents balance at
$b=\alpha(s+1)/(1+\alpha)$, where their common value is
$\alpha(s+1)/(1+\alpha)$. Its gate coefficient approaches $1/4$.
Section~\ref{sec:consistency-graph} proves the graph construction,
reductions, and frontier bound; Corollary~\ref{cor:gate-counting}
proves the uniform count-preserving algorithm and its bit complexity.

\paragraph{Status and scope.}
This is an asymptotic algorithm with exponential space, derived from
established width and tensor methods with explicit circuit-specific
accounting. The repository contains finite checks of supporting
constructions and a Rust implementation of the complete automatic
counting route in \texttt{rs/}. Its source-to-proof correspondence and
fixed-parameter resource argument are recorded in \texttt{rs/ALGORITHM.md}.
The retained structured timings illustrate a mechanism; they do not
establish the asymptotic bound empirically. Historical priority of the precise bound and its
best-known unrestricted-space status remain unestablished; the
statement-level comparison is in Section~\ref{sec:prior-work}.

\paragraph{From SAT to understanding nondeterminism.}
The proof first asks what changes when only witness inputs
are quantified and ordinary inputs remain symbolic. The same
consistency budget then gives a nonuniform circuit-size exponent
approaching $1/5$, alongside witness-coverage bounds, necessary
conditions for strong amplification, and examples separating a
complicated supplied circuit from an intrinsically hard function.

Returning to integer tables then gives the SAT bound and its consequences
for search, indexing, and exact generation in
Section~\ref{sec:algorithm-consequences}.
